\section{History as Environment Driven Recursive Self-Improvement}
\label{sec:method}

We consider a scientific discovery process that iteratively generates and
evaluates candidate solutions. At each step, an \emph{exploration policy}
decides how search effort should be allocated---which directions to open, which
existing branches to refine, and when to prune or stop. Rather than fixing this
exploration policy throughout discovery, our goal is to improve the exploration process
itself from the experience it has already generated.

Our key idea is to transform accumulated exploration history into a
\emph{replayable environment} for learning exploration policy. After each
discovery cycle, we retain the complete evaluated exploration trace rather than
only the best solution. These traces expose the same state--action interface as
live discovery, allowing alternative controllers to make exploration decisions
and observe their historically realized consequences without rerunning the
underlying candidate-generation process. A controller-development agent then
proposes and improves executable controllers through interaction with these
historical environments. The selected controller is deployed in the next live
discovery cycle, whose resulting experience expands the environment for
subsequent improvement.

Formally, let $H_t$ denote all exploration experience accumulated through
discovery cycle $t$. Our recursive improvement loop is
\begin{equation}
    H_t
    \;\longrightarrow\;
    \mathcal{E}(H_t)
    \;\longrightarrow\;
    \pi_{t+1}
    \;\longrightarrow\;
    \tau_{t+1}
    \;\longrightarrow\;
    H_{t+1}=H_t\cup\{\tau_{t+1}\},
    \label{eq:outer-loop}
\end{equation}
where $\mathcal{E}(H_t)$ is the historical environment constructed from
accumulated traces, $\pi_{t+1}$ is the learned exploration controller, and
$\tau_{t+1}$ is the new trace generated by deploying it. The framework thus
consists of three stages: constructing a historical environment, improving
exploration controllers within that environment, and recursively deploying the
improved controller to generate new experience.


\subsection{Constructing the Historical Environment}
\label{sec:historical-environment}

\paragraph{Exploration history.}
Each discovery cycle explores a branching space of candidate solutions. A root
node opens a new search direction from the task seed, while subsequent nodes
refine previous attempts along the same branch. For every evaluated node, we
record its branch structure, parent, candidate artifact, evaluator score,
validity, failure information, and task-specific diagnostics. We denote the
observation associated with node $v$ by
\begin{equation}
    o_v =
    \bigl(
        b,d,s_v,z_v,
        \Delta_v^{\mathrm{seed}},
        \Delta_v^{\mathrm{parent}}
    \bigr),
\end{equation}
where $(b,d)$ identifies the branch and depth, $s_v$ is the evaluator score,
$z_v$ contains evaluation status and diagnostics, and the two deltas measure
progress relative to the cycle seed and the parent node. Failed attempts are
retained because they provide evidence about unsuccessful directions and
potentially recoverable errors.

A completed discovery cycle produces a frozen exploration trace
$\tau=(V,E,O)$ together with its execution metadata. Importantly, we preserve
the explored structure and intermediate outcomes rather than only the final
best candidate. The cumulative history
\begin{equation}
    H_t = \{\tau_1,\ldots,\tau_t\}
\end{equation}
therefore captures how discovery unfolded across previous cycles, including
successful trajectories, regressions, failures, and unexplored alternatives
within the recorded search structure.

\paragraph{From trace to environment.}
We turn these frozen traces into a stateful decision environment. At decision
step $k$, a controller observes only the revealed prefix
\begin{equation}
    P_k = \{(v,o_v): v \text{ has been probed}\}.
\end{equation}
Given $P_k$, the legal action set $\mathcal{A}(P_k)$ contains unopened roots and
the next unrevealed child of each opened branch whose parent has already been
revealed. The controller selects a batch
\begin{equation}
    B_k \subseteq \mathcal{A}(P_k),
    \qquad |B_k|\leq W,
\end{equation}
where $W$ is the available worker capacity. This interface captures both
sequential exploration decisions and parallel allocation of search effort.

The key design is that \emph{live discovery and historical replay share the
same controller interface}. In live mode, probing a node launches the
candidate-generation agent and evaluator and returns the resulting observation.
In replay mode, the identical action reveals the observation previously
recorded for that node. State construction, legal actions, controller inputs,
probe accounting, and parallel execution accounting are otherwise shared.
Consequently, a controller developed through replay is directly executable in
live discovery without conversion from a simulator-specific policy.

The historical environment is deliberately restricted to observed support.
Replay can test how a controller would select, order, batch, prune, or stop
among previously evaluated nodes, but it never predicts the outcome of a
candidate that was not generated in the original trace. We therefore reuse the
\emph{decision structure} and realized consequences of previous exploration
without fabricating counterfactual candidate outcomes.


\subsection{Improving Exploration Controllers}
\label{sec:controller-improvement}

\paragraph{Executable controllers.}
An exploration controller $\pi$ is an executable program that maps the
currently revealed exploration state to future search decisions:
\begin{equation}
    \pi:
    \bigl(P_k,\mathcal{A}(P_k)\bigr)
    \longrightarrow B_k .
\end{equation}
The controller operates at two resolutions. Before a discovery cycle, it
selects a structural exploration envelope
\begin{equation}
    g=(M,D),
\end{equation}
where $M$ determines the available search width and $D$ the maximum refinement
depth. Within that envelope, the controller adaptively decides which roots to
open, which branches to refine or recover, and when further exploration is no
longer worthwhile.

Controllers are constrained to be deterministic and prefix-only. They may use
observed branch trajectories, relative improvements, failure evidence, and
remaining legal search depth, but cannot access unrevealed outcomes, winning
node identities, or the final best score of a historical trace. These
restrictions ensure that replay exposes only information that would have been
available at the corresponding point in live discovery.

\paragraph{Agentic controller improvement.}
Given $H_t$, we instantiate the historical environment
$\mathcal{E}(H_t)$ and provide a controller-development agent with the
controller API, previous controller implementations, and feedback from earlier
proposal rounds. The agent iteratively proposes executable controller code,
evaluates it through replay, inspects its behavior, and revises the controller:
\begin{equation}
    \pi^{(j)}
    \xrightarrow{\;\mathcal{E}(H_t)\;}
    F^{(j)}
    \longrightarrow
    \pi^{(j+1)},
    \label{eq:controller-learning}
\end{equation}
where $F^{(j)}$ summarizes replay performance and decision-level execution
feedback. Such feedback exposes behaviors such as premature stopping,
over-pruning, excessive probing, poor recovery decisions, or insufficient use
of parallelism. The development agent can use these observations to modify the
controller implementation in subsequent proposal rounds.

This design separates the \emph{learner} from the \emph{learned object}: the
controller-development agent performs meta-level optimization, while the
resulting controller is a lightweight executable policy that governs live
exploration.

\paragraph{Replay objective.}
A useful exploration controller should recover strong outcomes while spending
as little search effort as possible. For historical trace $i$ and exploration
setting $\beta$, let $Q_{i,\beta}$ denote the number of probes selected by the
controller, $S_{i,\beta}$ its best revealed score, $G_i$ the best score
available anywhere in the trace, and $N_i$ the number of evaluated nodes.
We measure normalized attainment as
\begin{equation}
    a_{i,\beta}
    =
    \operatorname{clip}
    \left(
        \frac{S_{i,\beta}}{G_i},0,1
    \right).
\end{equation}
Here, $G_i$ is used only for offline evaluation and is never exposed to the
controller.

Sweeping a fixed set of exploration settings yields a probe--attainment
frontier for each trace. We summarize this frontier using its normalized area,
$\operatorname{AUC}_i(\pi)$, which rewards controllers that recover high-quality
historical outcomes without exhaustively probing the trace. Because probe
efficiency alone can favor highly sequential policies that are inefficient in
live discovery, we additionally penalize unnecessary sequential decision
rounds. The final controller objective is
\begin{equation}
    J(\pi)
    =
    \frac{1}{|\mathcal{T}|}
    \sum_{i\in\mathcal{T}}
    \operatorname{AUC}_i(\pi)
    -
    \lambda
    \frac{1}{|\mathcal{T}|}
    \sum_{i\in\mathcal{T}}
    P_i,
    \label{eq:controller-objective}
\end{equation}
where $P_i$ measures the controller's sequential-round cost relative to its
probe usage on trace $i$. We select the highest-scoring valid controller across
proposal rounds. Detailed definitions of the frontier metric, parallelism
penalty, and controller validation are provided in Appendix~\ref{app:controller}.


\subsection{Recursive Deployment}
\label{sec:recursive-deployment}

After controller development, the selected controller $\pi_{t+1}$ is deployed
directly in the next live discovery cycle through the same interface used
during historical replay. Its actions now launch real candidate-generation and
evaluation jobs rather than revealing stored outcomes. The resulting
exploration trace $\tau_{t+1}$---including successful candidates, failures,
branch trajectories, scores, and execution decisions---is frozen and appended
to the cumulative history:
\begin{equation}
    H_{t+1}
    =
    H_t \cup \{\tau_{t+1}\}.
\end{equation}

The enlarged history then defines a new historical environment
$\mathcal{E}(H_{t+1})$ for learning the next controller. Because the replay pool
is cumulative, later controllers are evaluated across environments generated
by multiple previous controllers rather than only the most recent trajectory.
Each deployed controller therefore changes the experience distribution from
which its successor is learned, closing the recursive loop
\begin{equation}
    H_t
    \rightarrow
    \mathcal{E}(H_t)
    \rightarrow
    \pi_{t+1}
    \rightarrow
    H_{t+1}
    \rightarrow
    \mathcal{E}(H_{t+1})
    \rightarrow
    \pi_{t+2}
    \rightarrow \cdots .
\end{equation}

In summary, \textbf{History as Environment} converts the by-product of
discovery into a reusable substrate for improving future discovery. Historical
exploration provides the environment, a controller-development agent learns
how search should be conducted within that environment, and live deployment
generates the experience from which the next exploration controller can be
improved.